\section{附录}

\subsection{支撑材料与复现环境}

最终提交文件统一存放于 \path{./exports}。本次不生成论文 PDF；最终 PDF 文件名约定为
\path{第2次第66组.pdf}，支撑材料存放于 \path{./exports/支撑材料}。

\path{exports/支撑材料/复现环境} 是不依赖本仓库其他目录的独立程序副本。副本中的程序内容与
仓库 \path{src} 中的对应源程序一致，仅将文件名改为中文，并保留了程序所需的相对目录层级。
从复现环境根目录运行下列命令即可重新生成相应的处理中间数据、结果表、图件和核验日志；程序
运行时会在该复现环境内部创建 \path{data/processed} 与 \path{outputs}，不读取或改写赛题原始数据。

\begin{itemize}
  \item 问题一 Python 程序：\verb|python src/q1/问题一_静水风载系泊系统求解器.py|；
  \item 问题一 MATLAB 程序：\verb|matlab -batch "run('src/q1/问题一_静力模型.m')"|；
  \item 问题二完整 MATLAB 程序：\verb|matlab -batch "run('src/q2/问题二_完整程序.m')"|；
  \item 问题二静力模型程序：\verb|matlab -batch "run('src/q2/问题二_静力模型.m')"|；
  \item 问题三 Python 程序：\verb|python src/q3/问题三_风流荷载工程计算.py|（需要 Matplotlib）。
\end{itemize}

MATLAB 程序按源程序注释使用 MATLAB R2025b；Python 问题一程序只使用标准库，问题三程序使用
标准库和 Matplotlib。程序、数据、输出图表、结果表与核验日志之间通过问题编号和中文文件名对应。

\subsection{可运行源程序清单}

\begin{itemize}
  \item \path{exports/支撑材料/复现环境/src/q1/问题一_静水风载系泊系统求解器.py}；
  \item \path{exports/支撑材料/复现环境/src/q1/问题一_静力模型.m}；
  \item \path{exports/支撑材料/复现环境/src/q2/问题二_完整程序.m}；
  \item \path{exports/支撑材料/复现环境/src/q2/问题二_静力模型.m}；
  \item \path{exports/支撑材料/复现环境/src/q3/问题三_风流荷载工程计算.py}。
\end{itemize}

\subsection{处理后数据清单}

以下文件是模型计算产生的中间数据或参数扫描数据；赛题提供的原始数据未复制到支撑材料中。

\begin{itemize}
  \item \path{exports/支撑材料/处理后数据/问题一_锚链形坐标_12米每秒.csv}；
  \item \path{exports/支撑材料/处理后数据/问题一_锚链形坐标_24米每秒.csv}；
  \item \path{exports/支撑材料/处理后数据/问题一_吃水解方程.csv}；
  \item \path{exports/支撑材料/处理后数据/问题一_水深与吃水曲线数据.csv}；
  \item \path{exports/支撑材料/处理后数据/问题二_锚链坐标_1200千克.csv}；
  \item \path{exports/支撑材料/处理后数据/问题二_锚链坐标_2250千克.csv}；
  \item \path{exports/支撑材料/处理后数据/问题二_配重质量扫描.csv}；
  \item \path{exports/支撑材料/处理后数据/问题三_锚链形状_16米.csv}；
  \item \path{exports/支撑材料/处理后数据/问题三_锚链形状_20米.csv}；
  \item \path{exports/支撑材料/处理后数据/问题三_水深与吃水曲线数据.csv}。
\end{itemize}

\subsection{结果图表清单}

\begin{itemize}
  \item 问题一：\path{exports/支撑材料/结果图表/问题一/问题一_锚链形状.pdf}、
    \path{exports/支撑材料/结果图表/问题一/问题一_锚链形状.png}、
    \path{exports/支撑材料/结果图表/问题一/问题一_锚链形状.svg}、
    \path{exports/支撑材料/结果图表/问题一/问题一_锚链形状.tex}；
    \path{exports/支撑材料/结果图表/问题一/问题一_水深与吃水曲线.pdf}、
    \path{exports/支撑材料/结果图表/问题一/问题一_水深与吃水曲线.png}、
    \path{exports/支撑材料/结果图表/问题一/问题一_水深与吃水曲线.svg}、
    \path{exports/支撑材料/结果图表/问题一/问题一_水深与吃水曲线.tex}；
  \item 问题二：\path{exports/支撑材料/结果图表/问题二/问题二_锚链形状对比.pdf}、
    \path{exports/支撑材料/结果图表/问题二/问题二_锚链形状对比.png}、
    \path{exports/支撑材料/结果图表/问题二/问题二_配重约束关系.pdf}、
    \path{exports/支撑材料/结果图表/问题二/问题二_配重约束关系.png}；
  \item 问题三：\path{exports/支撑材料/结果图表/问题三/问题三_16米与20米锚链形状.pdf}、
    \path{exports/支撑材料/结果图表/问题三/问题三_16米与20米锚链形状.png}、
    \path{exports/支撑材料/结果图表/问题三/问题三_16米与20米锚链形状.svg}、
    \path{exports/支撑材料/结果图表/问题三/问题三_水深与吃水曲线.pdf}、
    \path{exports/支撑材料/结果图表/问题三/问题三_水深与吃水曲线.png}、
    \path{exports/支撑材料/结果图表/问题三/问题三_水深与吃水曲线.svg}。
\end{itemize}

\subsection{结果表与核验日志清单}

\begin{itemize}
  \item 问题一：\path{exports/支撑材料/结果表与日志/问题一/问题一_静力结果表.csv}、
    \path{exports/支撑材料/结果表与日志/问题一/问题一_静力结果表.tex}、
    \path{exports/支撑材料/结果表与日志/问题一/问题一_静力核验日志.txt}；
  \item 问题二：\path{exports/支撑材料/结果表与日志/问题二/问题二_配重扫描.csv}、
    \path{exports/支撑材料/结果表与日志/问题二/问题二_配重扫描.tex}、
    \path{exports/支撑材料/结果表与日志/问题二/问题二_静力结果表.csv}、
    \path{exports/支撑材料/结果表与日志/问题二/问题二_静力结果表.tex}、
    \path{exports/支撑材料/结果表与日志/问题二/问题二_静力核验日志.txt}、
    \path{exports/支撑材料/结果表与日志/问题二/问题二_结果摘要.txt}；
  \item 问题三：\path{exports/支撑材料/结果表与日志/问题三/问题三_重物球质量影响.csv}、
    \path{exports/支撑材料/结果表与日志/问题三/问题三_重物球质量影响.tex}、
    \path{exports/支撑材料/结果表与日志/问题三/问题三_流速与锚链模型关系.csv}、
    \path{exports/支撑材料/结果表与日志/问题三/问题三_流速与锚链模型关系.tex}、
    \path{exports/支撑材料/结果表与日志/问题三/问题三_流速影响.csv}、
    \path{exports/支撑材料/结果表与日志/问题三/问题三_流速影响.tex}、
    \path{exports/支撑材料/结果表与日志/问题三/问题三_水深影响.csv}、
    \path{exports/支撑材料/结果表与日志/问题三/问题三_水深影响.tex}、
    \path{exports/支撑材料/结果表与日志/问题三/问题三_链节数影响.csv}、
    \path{exports/支撑材料/结果表与日志/问题三/问题三_链节数影响.tex}、
    \path{exports/支撑材料/结果表与日志/问题三/问题三_静力链模型关系.csv}、
    \path{exports/支撑材料/结果表与日志/问题三/问题三_静力链模型关系.tex}、
    \path{exports/支撑材料/结果表与日志/问题三/问题三_风速影响.csv}、
    \path{exports/支撑材料/结果表与日志/问题三/问题三_风速影响.tex}、
    \path{exports/支撑材料/结果表与日志/问题三/问题三_工程核验日志.txt}。
\end{itemize}

\subsection{模型说明材料}

\begin{itemize}
  \item \path{exports/支撑材料/模型说明/问题一_模型卡.md}；
  \item \path{exports/支撑材料/模型说明/问题一_数据说明.md}。
\end{itemize}

上述文件均为支撑模型、结果和结论所需的程序、处理中间数据、图表、结果表、核验日志或模型
说明材料；未复制赛题提供的原始数据，也未另行使用未列入清单的外部数据资料。
